\documentclass[11pt]{article}
\usepackage[a4paper,margin=1in]{geometry}
\usepackage{amsmath,amssymb,amsthm,mathtools}
\usepackage{booktabs,tabularx,array}
\usepackage{enumitem}
\usepackage[T1]{fontenc}
\usepackage{lmodern}
\usepackage{microtype}
\usepackage[hidelinks]{hyperref}
\usepackage{xurl}
\newcolumntype{Y}{>{\raggedright\arraybackslash}X}
\newtheorem{definition}{Definition}
\newtheorem{proposition}{Proposition}
\newtheorem{remark}{Remark}
\newcommand{\RAKL}{\textsc{Orion}}
\newcommand{\code}[1]{\texttt{#1}}

\title{Structural Learning Mechanics:\\
Instrument Admissibility for Equal-Budget Allocation Comparisons\\
\large A negative-results and metrology report}
\author{Sze Chun Yiu\\Stockholm University\\\texttt{sze-chun.yiu@fysik.su.se}}
\date{August 2026}

\begin{document}
\sloppy
\maketitle

\begin{abstract}
An experiment can be amply powered by every conventional criterion and still be structurally unable to answer its own registered question. This paper reports that failure mode concretely, names the metrology defect, and contributes the method that detects it before a confirmatory budget is spent. The empirical core is a preserved negative: in a frozen model-free development-stress instrument, a learner-conditioned adaptive allocation policy lost to a static structural parent under equal training budget ($E-D=-0.0166$, bootstrap 95\% CI $[-0.0174,-0.0158]$), with the loss attributed by a frozen diagnostic protocol to the allocation stage --- guard-rail budget capture plus within-round concentration --- and with the protocol's own pre-registered attribution prediction \emph{falsified} and preserved. The decisive finding sits one level deeper: a rigorous upper bound on the advantage \emph{any} equal-budget allocation policy could achieve over the static parent in that instrument is $+0.0246$, a factor $2.03$ below the instrument's own frozen $0.05$ promotion gate. The instrument could not have produced a positive verdict for any policy, however good --- and this is invisible to conventional power analysis, because the bootstrap confidence intervals on the primary contrast are ${\approx}0.0016$ wide. Power analysis bounds sampling noise for an assumed effect size; it is silent on whether the instrument can generate the effect at all. We formalize the repair as an \emph{oracle-ceiling admissibility gate}: before executing an equal-budget allocation comparison, compute bounds on the achievable ceiling and require $\text{ceiling}\ge\kappa\cdot\mathrm{MDE}$ with $\kappa$ frozen in advance; only a rigorous upper bound may license \code{INADMISSIBLE}, only a constructive achievability bound may license \code{ADMISSIBLE}, and an uncomputable oracle fails closed to \code{CANNOT\_CHECK}. The gate passes a five-case known-answer fresh-assurance battery with pre-frozen expected verdicts and a per-condition black-box falsifiability audit. A demonstration closes the loop: repairing the diagnosed instrument defect (world-divergent initial learner state) yields a licensed instrument in which a marginal-gain challenger beats the static parent on every frozen gate ($F-D=+0.0760$) --- and in which the original ``defective'' v1 policy \emph{also} wins ($+0.0747$), falsifying a second registered prediction and showing the original negative was policy$\times$instrument-conditional. Both falsified predictions are preserved verbatim. The paper makes no positive claim about real-model training; the original standalone positive-claims Paper~IV remains unauthorized, and this report does not reverse that decision.
\end{abstract}

\begin{center}
\fbox{\parbox{0.94\textwidth}{\small
\textbf{Status: rescoped claims object.} An earlier version of this document was a design-and-protocol preprint for a prospective positive adaptive-allocation claim. That standalone positive paper was, and remains, unauthorized under decision gate \#462; this rescope does not reverse or relitigate that decision --- it is a \emph{different claims object}: a negative-results and metrology report whose every quantitative statement is read from preserved, frozen receipts in the accompanying research artifact. The preserved negatives are not reinterpreted as positives anywhere in this document. All analysis is same-context (produced within the same research programme, with language-model assistance); none of it is independent review.
}}
\end{center}

\section{The gap, the residual question, and what this paper now claims}
\label{sec:gap}

A learner does not benefit equally from every training example, and the modern data-selection literature acts on this at scale: skill-based frameworks sample adaptively along learned skill dependencies \cite{chen2023skillit}; skill-graph selection curates by structural role \cite{li2025mass}; model-aware influence methods track which data most improve the current parameters \cite{yu2024mates}; token-level methods select within an example \cite{lin2024rho1}; optimizer-aware influence methods target instruction tuning \cite{xia2024less}; de-duplication and diversification beat random one-pass exposure \cite{tirumala2023d4,abbas2023semdedup}; mixture reweighting optimizes domain proportions \cite{xie2023doremi}; and skill-targeted adaptive training profiles a learner's missing skills after saturation \cite{he2026stat}. Compositional-generalization work shows that atomic-skill examples do not automatically teach composition \cite{wei2026steps}; structural information is learned and used at inference \cite{chen2026structural}; latent skills route and compose across adaptation \cite{ling2026rex}; and spaced, varied retrieval improves retention and transfer, not merely acquisition \cite{butler2010}.

The research programme behind this paper registered a narrow residual hypothesis over that prior art: that allocation conditioned on an explicit, directional, quantity-of-interest-scoped \emph{structural} coordinate with a learner-specific \emph{vector} mastery state beats strong static structural curation under equal budget and noncompensatory safety constraints. That hypothesis is \emph{not} what this paper claims. What this paper now claims is what its receipts support:

\begin{enumerate}[leftmargin=1.5em,itemsep=2pt]
\item \textbf{A metrology method (primary contribution).} An equal-budget allocation comparison is uninformative about its registered estimand when the achievable oracle ceiling of the instrument lies below its registered minimum detectable effect, and this failure mode is invisible to conventional power analysis. We give a concrete admissibility gate --- $\text{ceiling}\ge\kappa\cdot\mathrm{MDE}$, $\kappa$ frozen before ceiling access, with direction-of-license rules and fail-closed semantics --- demonstrated on exactly one preserved instrument and validated on a fresh known-answer battery (\S\ref{sec:admissibility}--\S\ref{sec:gatevalidation}).
\item \textbf{An attributed negative (empirical core).} The adaptive-v1 policy lost to its static parent in the frozen development-stress instrument; the loss decomposes into guard-rail budget capture and within-round concentration; and the diagnostic protocol's own pre-registered attribution prediction failed and is preserved as failed (\S\ref{sec:negative}).
\item \textbf{A closed-loop demonstration (secondary).} Repairing the diagnosed instrument defect produces a licensed instrument in which a marginal-gain challenger clears every frozen gate against the static parent --- and in which the v1 policy also wins, falsifying a second registered prediction and relocating the original negative from ``policy defect'' to ``policy$\times$instrument interaction'' (\S\ref{sec:demonstration}).
\end{enumerate}

Everything else --- the structural-coordinate mechanics, the preregistered Phase~0--4 protocol, the frozen 7B confirmatory design --- appears here as context and preserved history, not as claims. No real-model result exists, and none is asserted.

\section{Structural coordinates and the mechanics under test}
\label{sec:mechanics}

The mechanics are retained from the design phase because the negative results below are unintelligible without them; they are architecture, not findings.

Let $R_t$ denote the external research substrate at time $t$. The framework exposes a scientific-authority projection $K_t=\pi_{\mathrm{epi}}(R_t)$ and a search projection $Q_t=\pi_{\mathrm{search}}(R_t)$; the training extension adds $L_t=\pi_{\mathrm{train}}(R_t,\theta_t)$, the only projection that depends on model parameters $\theta_t$. The binding invariant is that training utility, search rank and scientific authority never share semantics.

\begin{definition}[Structural mastery vector]
\label{def:mastery}
For a registered structural object $s$ and parameters $\theta_t$, the learner-conditioned mastery state is the vector
\[
M_t(s)=\bigl(m_{\mathrm{principle}},\,m_{\mathrm{composition}},\,m_{\mathrm{boundary}},\,m_{\mathrm{representation}},\,m_{\mathrm{transfer}},\,m_{\mathrm{retention}}\bigr),
\]
each coordinate estimated by a frozen probe family rather than model introspection.
\end{definition}

Three architectural propositions from the design phase remain load-bearing for reading the results: learner-conditioned saturation is not a fixed point of the training operator (a saturation receipt is staleable and checkpoint-bound); a scalar mastery summary is representationally inadequate for the allocation decision (two states with equal mean demand different allocations); and retention/transfer safety must be noncompensatory constraints, since any finite additive retention weight admits schedules that buy the primary objective with unbounded forgetting \cite{tirumala2023d4,butler2010}. The production allocator implements these as hard gates; the negative below is \emph{not} a failure of these constraints but of the allocation policy operating inside them.

\paragraph{The estimand.}
The registered causal object is the contrast $E-D$ at equal total budget between \textbf{E}, adaptive structural-mastery allocation, and \textbf{D}, \emph{static} structural curation over the identical substrate, alongside \textbf{A} uniform random, \textbf{B} semantic diversity, and \textbf{C}, the strongest scalar loss-aware learner-conditioned parent. $E-D$ isolates the learner-conditioning residual; $E-A$ and $E-B$ conflate it with curation value; $E-C$ tests residual value beyond generic learner-conditioning.

\section{Preserved instrument history I: the Phase-0/1 retraction}
\label{sec:phase1-result}

The programme's first instrument lesson predates the present one and is preserved because the two share a defect class. The Phase-0/1 exposure instrument (frozen packet, subject hash \texttt{fce2bb17\ldots}) was executed on the frozen Qwen2.5 instruction-tuned ladder ($0.5$B--$7$B, A100, LoRA per exposure count, six mastery coordinates on held-out cases) --- and then \textbf{retracted as an instrument artifact} on root-cause analysis. Two independent defects were found: (1) a \emph{degenerate generator} exposing only two unique rendered inputs per structural family, so no structural rule could be induced and an edge-count heuristic explained the apparent learning; (2) a \emph{masked-answer-token training bug}: byte-pair encoding merged the separating space into the answer token, the masking boundary fell after the gold token, and no gradient ever reached the VALID/INVALID decision --- every reported outcome was the untrained base model's prior. Either defect alone voids the run, so it identifies neither a residual nor a capability floor. The v1 terminals are preserved as negative instrument history supporting no claim in either direction:
\begin{gather*}
\texttt{V1\_PHASE1}=\texttt{INSTRUMENT\_GENERATOR\_DEFECT},\quad
\texttt{MECHANISM\_SIGNAL}=\texttt{CANNOT\_CHECK}.
\end{gather*}
A corrected v2 instrument (token-id concatenation; varied, length-matched, disjoint train/probe instances; a learnability positive-control gate) restored genuine generalising learning, and its A100 re-run produced one real signal on the corrected ladder: 7B \code{state\_reachability} $\to$ \code{MECHANISM\_SIGNAL\_PRESENT} (principle mastered at exposure 2, late same-structure gain $0$, unsaturated-coordinate late gain $+0.25$), in exactly 1 of 12 model$\times$family cells. That signal motivated the allocator whose failure is this paper's empirical core. The pattern to hold onto: \emph{an instrument that cannot produce the signal it was built to test, while looking healthy}. The Phase-0/1 v1 instrument failed this way at the level of gradient flow; the development-stress instrument below fails the same way at the level of achievable effect size.

\section{The empirical core: an attributed equal-budget negative}
\label{sec:negative}

\subsection{The preserved development negative}
A model-free development-stress panel was frozen before outcomes (seed $202608140601$; six worlds spanning early principle saturation, composition/boundary/representation/transfer lag, and retention sensitivity; $384$ world-replicates; equal $48$-example budget per arm; balanced mastery and hard-safety-minimum endpoints; promotion gate $E-D\ge 0.05$). The adaptive-v1 policy --- one PRINCIPLE floor slot per round, principle-until-$0.90$, hard retention repair, otherwise all remaining slots to the weakest measured coordinate --- \textbf{lost}:
\[
E-D\ \text{(balanced mastery)} = -0.0166\ [-0.0174,-0.0158],\qquad
E-D\ \text{(hard-safety min)} = -0.0503\ [-0.0578,-0.0431],
\]
preserved as \code{DEVELOPMENT\_NEGATIVE\_ADAPTIVE\_OVERCONCENTRATES\_\allowbreak STATIC\_PARENT\_RETAINS\_DEFAULT}, with active-default authority revoked from adaptive scheduling: the static structural parent remains the governed production policy. This outcome was produced by the closure lane of the research programme \emph{before} any of the admissibility analysis below existed; the analysis inherits the outcome, it did not generate it.

\subsection{Frozen attribution, with its own falsified prediction}
A follow-up diagnostic protocol was frozen before outcome access with four numbered predictions and falsifiers, then executed with world dynamics and RNG streams imported unchanged from the frozen parent runner. Findings, all read from receipts:

\begin{itemize}[leftmargin=1.4em,itemsep=2pt]
\item \textbf{Guard-rail budget capture.} v1 spends $13$ of $48$ examples on PRINCIPLE versus the parent's $8$ --- in every world --- because the principle threshold rule and the floor slot are budget-consuming \emph{targets} while the dynamics simultaneously erode PRINCIPLE. Per-coordinate $E-D$ mastery: PRINCIPLE $+0.0387$, every other coordinate negative. Capping PRINCIPLE at the parent's $8$ (arm X1) drives the allocation's $L_1$ distance from uniform from $10.33$ to $1.31$ and recovers $35\%$ of the gap.
\item \textbf{Within-round concentration.} X1 reaches near-uniform budget yet still loses $-0.0108$: committing seven slots to one coordinate inside a round wastes budget against per-coordinate saturating gains.
\item \textbf{Falsified prediction P1, preserved.} The protocol's pre-registered prediction P1 stated that the principle cap closes ${\ge}60\%$ of the gap ($\Delta\ge 0.010$). Observed: $\Delta=+0.0058$, $35\%$ of the gap. \textbf{P1 is recorded as FALSIFIED.} No threshold was moved after outcome access; the attribution terminal is \code{ATTRIBUTION\_MIXED}, not the single-lever story P1 predicted.
\item \textbf{A confounded arm, declared uninformative.} Arm X2, designed to isolate concentration, instead \emph{increased} budget distortion ($L_1$ $10.33\to16.00$) because spreading slots re-selects the argmin every round; its deficit is fully accounted for by budget distortion, and it is used as evidence in \emph{neither} direction.
\item \textbf{A corrected root-cause reading.} The original receipt's narrative (``commits all non-repetition slots to one weakest coordinate'') is refuted \emph{as a budget-level story} by the realized counts, which are near-uniform outside PRINCIPLE; the receipt file itself is immutable and unedited, and only its interpretation is corrected.
\end{itemize}

Both v1 defects --- argmin mastery \emph{level} where the thesis concerns the mastery \emph{derivative}, and whole-round single-target blocks --- are present in the promoted production allocator, so the diagnosis transfers beyond the simulator.

\section{The primary contribution: oracle-ceiling instrument admissibility}
\label{sec:admissibility}

\subsection{The decisive measurement}
The attribution's deepest finding is not about the policy at all. Three bounds were computed on the advantage \emph{any} equal-budget allocation policy could achieve over the static parent in this instrument:

\begin{center}
\small
\begin{tabular}{@{}llll@{}}
\toprule
tier & object & direction & mean advantage over static \\
\midrule
1 & greedy oracle \emph{policy} (true-parameter access) & lower bound & $+0.00148$ \\
2 & best constructive allocation found (local search) & lower bound & $+0.00446$ \\
3 & harm-free relaxation, solved exactly & \textbf{upper bound} & $\mathbf{+0.02457}$ \\
\bottomrule
\end{tabular}
\end{center}

The tier-3 bound is rigorous for this instrument: every harm term in the dynamics is non-negative and subtracted, the gain step $m\mapsto m(1-r)+r$ is monotone increasing in $m$, so the harm-free trajectory dominates every with-harm trajectory pointwise; maximizing it over integer count vectors is a concave separable allocation solved exactly by greedy water-filling. Hence \emph{no equal-budget allocation policy, however optimal, can reach this instrument's own registered $0.05$ gate} --- the achievable ceiling sits a factor $2.03$ below it.

\paragraph{Correction record, preserved.} A first-pass analysis cited the greedy oracle's $+0.00148$ as ``the ceiling, $34\times$ below the gate.'' That inference was invalid --- a greedy oracle is a \emph{policy}, so its score is only a \emph{lower} bound on the ceiling, and it understated the constructively achievable advantage by a factor of three. The correction strictly weakens the headline margin ($34\times\to 2.03\times$) and was made before the gate was implemented or assessed. The tier-3 relaxation is loose by construction (tier-2/tier-3 spread ${\approx}5.5\times$), so the margin is stated as a factor of two on a loose bound, not as airtight.

\subsection{Invisible to power analysis}
The bootstrap 95\% confidence intervals on the primary contrast are ${\approx}0.0016$ wide --- by conventional power analysis this experiment is amply powered, and no pre-registration checklist item flags it. The defect is orthogonal to sampling noise: classical power analysis bounds the noise around an \emph{assumed} effect size; it is silent on whether the instrument's dynamics can generate an effect of that size at all. Three structural causes were identified and verified in the frozen runner: round-0 learner state is effectively world-independent, so state information carries almost nothing a static arm does not exploit; gains are concave and separable while harms are linear and coordinate-independent, making equal spread near-optimal by construction; and six rounds against five non-principle coordinates degenerate argmin cycling into round-robin. Consequently the preserved negative is valid \emph{only} as ``this v1 policy in this instrument'': the instrument could not have adjudicated the registered estimand in either direction, for any policy.

\subsection{The admissibility gate}
\label{sec:gatedef}
The repair, frozen as a sealed mechanic packet before implementation, is a pre-execution gate:

\begin{definition}[Oracle-ceiling admissibility]
\label{def:gate}
Let an equal-budget allocation-comparison instrument register a primary contrast with minimum detectable effect $\mathrm{MDE}>0$, and let $\kappa>0$ be frozen before any ceiling value is computed. Given bounds on the instrument's achievable ceiling (the supremum over admissible policies of the mean primary-contrast advantage over the reference parent), the gate returns:
\code{INADMISSIBLE} iff a rigorous \emph{upper} bound is $<\kappa\cdot\mathrm{MDE}$;
\code{ADMISSIBLE} iff a constructive \emph{achievability} (lower or exact) bound is $\ge\kappa\cdot\mathrm{MDE}$;
and \code{CANNOT\_CHECK} otherwise --- including whenever the oracle is not computable under the instrument's own generative model, the budget constraint is unverified, or the supplied bounds are missing, inconsistent, or too loose to decide.
\end{definition}

The direction-of-license rules are the content of the mechanic: a policy score can prove an instrument \emph{has} headroom, never that it lacks it; a relaxation optimum can prove it \emph{lacks} headroom, never that any policy attains it. An \code{INADMISSIBLE} verdict may never be upgraded by access to arm outcomes; $\kappa$ and the MDE may never be revised after ceiling access; the gate decides only whether a comparison is worth executing and grants no authority of any kind. In the accompanying artifact, $\kappa=1.2$ is frozen with a declared rationale (an instrument whose best achievable effect sits exactly at its MDE gives the best possible policy a coin flip at its own gate), together with the disclosure that for the reference instrument the verdict is $\kappa$-\emph{insensitive} --- \code{INADMISSIBLE} holds for every $\kappa>0.4914$ --- so the choice of $\kappa$ is not outcome-tuned there.

\subsection{Relation to prior art, and the residual claimed}
Every ingredient is established: oracle/skyline baselines in retrieval and recommendation, regret-versus-oracle formulations in online allocation, ceiling analysis in pipeline debugging, and --- structurally closest --- the assay \emph{positive control}: an instrument that cannot register a known-positive signal is inadmissible. Classical power analysis and MDE pre-registration are the strongest incumbents and are what the frozen protocols here already used. The residual claimed is narrow: the conversion of an oracle/skyline computation into a \emph{frozen pre-execution admissibility decision bound to the registered MDE}, with fail-closed \code{CANNOT\_CHECK} when the oracle is not computable, and with the direction-of-license asymmetry made explicit. Demonstrated on exactly one preserved instrument; nothing more is claimed.

\section{Validating the validator}
\label{sec:gatevalidation}

A gate that cannot fail is worse than no gate, so the gate itself was audited before any confirmatory use.

\paragraph{Known-answer and hostile worlds.} The implementation carries one test per sealed-packet counterexample: exact zero headroom $\to$ \code{INADMISSIBLE}; genuine headroom $\to$ \code{ADMISSIBLE}; uncomputable oracle with a huge \emph{claimed} lower bound $\to$ \code{CANNOT\_CHECK}, never \code{ADMISSIBLE}; the observed tight-CI/low-ceiling case $\to$ \code{INADMISSIBLE} regardless of CI width; an exact ceiling in the open interval $(\mathrm{MDE},\kappa\cdot\mathrm{MDE})$ $\to$ \code{INADMISSIBLE} under the frozen $\kappa$, the no-post-hoc-rescue case; loose straddling bounds $\to$ \code{CANNOT\_CHECK}; a tiny lower bound alone can never license \code{INADMISSIBLE}; a huge upper bound alone can never license \code{ADMISSIBLE}.

\paragraph{Fresh assurance with pre-frozen verdicts.} Five fresh cases were constructed at a seed disjoint from every development artifact, with expected verdicts frozen \emph{before} the runner existed: all five verdicts matched, including an \code{ADMISSIBLE} licensed by a lower bound alone in a nonseparable-synergy instrument where no monotone relaxation exists (the gate works with no upper bound at all), and an exact-enumeration instrument swept deterministically until its true ceiling landed in $(\mathrm{MDE},\kappa\cdot\mathrm{MDE})$.

\paragraph{Black-box falsifiability audit, control first.} Each condition was probed by evidence perturbations with the expected \code{SENSITIVE}/\code{INSENSITIVE} outcome of every probe declared before execution --- including asserted \emph{no-alarm} rows, e.g.\ the \code{CANNOT\_CHECK} of an uncomputable oracle must be \emph{insensitive} to arbitrary inflation of the claimed bounds. All conditions pass. The audit's own first run produced two false alarms --- from the \emph{probe harness}, not the gate: independent per-bound scale factors occasionally produced lower$>$upper evidence, which the gate correctly refused as \code{CANNOT\_CHECK}. The probe was repaired to a shared draw; the gate was not changed; the correction is preserved in the runner.

\paragraph{Same-system ablation.} Replacing the oracle upper bound with the strongest implementable policy scores flips the reference verdict \code{INADMISSIBLE}$\to$\code{CANNOT\_CHECK} in $32/32$ trials: the upper-bound component carries the mechanic, and without it the gate cannot reduce to strongest-parent reporting.

\section{Closing the loop: a licensed instrument and an honest positive}
\label{sec:demonstration}

\subsection{Instrument repair and licensing}
The diagnosed blocking cause --- world-independent initial learner state --- was repaired in a successor instrument: each world assigns a distinct \emph{deficient coordinate pair} (initial mastery $0.10$ against near-saturated $0.94$ elsewhere), so which coordinates are weak is world-specific and observable only through the learner-state projection. Dynamics, total budget ($48$), rounds and batch size are imported unchanged from the frozen parent runner; PRINCIPLE and RETENTION each appear among the deficient pairs so that neither the v1 guard rails nor their removal is uniformly favoured. The gate licensed the first configuration tried --- \code{ADMISSIBLE}: constructive lower bound $+0.0800\ge\kappa\cdot\mathrm{MDE}=0.06$, upper bound $+0.0974$, a tier spread of $1.22\times$ against the reference instrument's $5.5\times$ --- and the design log contains exactly one entry.

The gate then did fresh pre-outcome work: a pre-declared rule required every registered contrast gate, not only the primary one, to be ceiling-licensed. The constructive ceiling over the scalar loss-aware parent C is only $+0.0046$ --- far below $\kappa\cdot 0.02$ --- so the v1-style ``beat C by $0.02$'' gate would have been a \emph{second} structurally unreachable gate; it was bound sign-only (CI excluding zero) before any arm outcome, and the material-margin question against C is recorded as unresolvable by this instrument.

\subsection{The challenger and the frozen result}
The queued policy lever was implemented as a versioned challenger, never mutating the production allocator: per-slot believed-marginal-gain water-filling ($\mathrm{argmax}\ \hat r_c(1-m_c)$ against a believed state that saturates as slots are spent, $\hat r$ estimated online from observed increments), with guard rails demoted from budget-consuming targets to constraints. All thresholds, seeds, falsifiers and terminal branches were frozen before outcome access; the assurance seed was consumed once; no algorithm change occurred between freeze and assurance. Result, read from the receipt:
\[
F-D = +0.0760\ [0.0755, 0.0764],\quad\text{positive in }6/6\text{ worlds},\quad
\text{safety } F-D = +0.157,
\]
clearing every frozen gate ($F$ extracts ${\approx}95\%$ of the licensed constructive ceiling), with $F-C=+0.00025$ positive under the sign-only gate --- at the scale of the $0.0046$ C-ceiling, exactly as the licensing analysis anticipated.

\subsection{A second falsified prediction, preserved}
The frozen diagnostic prediction P6 stated that the v1 policy would \emph{remain negative} in a licensed instrument --- that the original loss reflected an intrinsic policy defect. \textbf{P6 is falsified}: $E(\text{v1})-D=+0.0747$ in the licensed instrument. The honest reading, recorded in the receipts: the preserved parent negative was \emph{policy$\times$instrument-conditional}; in a state-divergent instrument, deficit-chasing of either kind (level- or gain-driven) beats static allocation, and the specific marginal-gain-versus-level contrast is small ($F-E\approx+0.0014$, both endpoints in the challenger's favour). The strongest supported claim is therefore \emph{adaptive-versus-static in licensed instruments}; the mechanism-level claim about the mastery derivative remains open. The parent negative itself is preserved verbatim and remains correctly scoped: static allocation was, and remains, the right production default on the evidence that existed.

\section{Preserved protocol context: the frozen confirmatory path}
\label{sec:protocol}

The preregistered Phase~0--4 programme and its governing decision records (\#461, \#466, \#467, \#468) remain frozen history: the exposure-ladder design, the matched five-arm Phase-2 comparison at the exact pinned Qwen2.5-7B revision with $\mathrm{MDE}=0.05$, disjoint selection probes, a $384$-case assurance panel, hard harm bounds and full overhead accounting. None of it has been executed at 7B, and nothing here substitutes for it. The present report adds one advisory, now backed by a formal verdict rather than a suspicion: the frozen 7B protocol registers $\mathrm{MDE}=0.05$ with no ceiling calibration, and whether that instrument has $0.05$ of state-conditioning headroom is \code{CANNOT\_CHECK} until probed --- the recommended cheap ceiling probe should precede the confirmatory spend. The frozen protocol files are untouched.

\section{Threats to validity}
\label{sec:threats}

\paragraph{The ceiling bound may be loose or wrong for other dynamics.} The tier-3 argument is instrument-specific: it requires non-negative subtractive harms and a monotone gain map. Nonseparable dynamics void the water-filling optimality argument --- in that case the gate must (and, in the fresh-assurance battery, does) either license via a constructive bound or fail closed. An \code{INADMISSIBLE} verdict inherits the rigor of the relaxation it stands on, and the reference verdict's $2.03\times$ margin sits on a bound that is loose by a factor of ${\approx}5.5$ against the best construction.

\paragraph{Model-free scope.} Every quantitative result is from model-free development-stress simulators with known generative parameters. That is precisely the regime the gate targets (the oracle is computable); for real-model instruments the oracle is not computable and the gate correctly returns \code{CANNOT\_CHECK} rather than a verdict. Nothing here transfers to real training runs except the method and the advisory.

\paragraph{The licensed instrument was designed to have headroom.} Instrument v2 was constructed so that state information is decision-relevant --- the analogue of spiking a positive control into an assay. The licensing ceiling is policy-agnostic and was computed before any arm outcome, and the sign-only C-gate shows the design was not tuned to flatter the challenger (its residual over C is three orders of magnitude below its residual over D). Still, a licensed instrument demonstrates that the comparison \emph{can} discriminate, not that its worlds are representative of any real training regime.

\paragraph{Same-context analysis.} The attribution, the gate, the licensing and the challenger were produced within one research programme (with the development negative produced by an earlier lane than the analysis). Same-context critique, however disciplined, is not independent review, and no claim here has received any.

\paragraph{Multiplicity across the revival lane.} The revival executed one gate battery, one licensing computation and one assurance run, with expected outcomes frozen before each; the two falsified predictions (P1, P6) are reported alongside the confirmed ones rather than being absorbed. Nonetheless the lane as a whole was run once, by one team, on one family of simulators.

\section{Relationship to the rest of the series}
\label{sec:series}

\textbf{Paper~I} contributes the product-order, no-scalarization and freeze-before-outcome discipline this report exercises. \textbf{Paper~II} contributes the directional structural objects the training projection annotates. \textbf{Paper~III} (method-evolution) is the closest sibling and the default owner of any conceptual residue of the adaptive-allocation programme; its own negative two-arm training precursor \cite{yiu2026verified} is not evidence for or against anything here. \textbf{Paper~V} contributes the routing-only-projection precedent and the discipline that proposal fluency never becomes authority. The admissibility mechanic itself is a candidate reusable tool for any sibling that runs equal-budget comparisons against registered MDEs; its consolidation into the shared framework is governed by the usual promotion gates, not by this paper.

\section{Scope, authorization, and what this paper does not do}
\label{sec:preconditions}

\begin{enumerate}[leftmargin=1.5em,itemsep=2pt]
\item \textbf{This rescope does not reverse \#462.} The original standalone positive-claims Paper~IV --- learner-conditioned structural allocation as a demonstrated mechanism --- was rejected at gate \#462 and remains rejected. Zero of that gate's justification conditions are met: the real-model exposure signal exists in 1 of 12 cells, the 7B Phase-2 comparison is unexecuted, and no cross-family law exists. This document is a different claims object (negative results and metrology) and claims nothing the gate governs.
\item \textbf{No training-policy authority moves.} The static structural parent remains the governed production default. The challenger's licensed-instrument win is development evidence for continuing its lane; activation would require the external \code{ADAPTIVE\_RESIDUAL\_SUPPORTED} path with fresh assurance, strongest-parent residual, hard-harm and full-overhead gates --- none of which this report touches.
\item \textbf{Negative history is immutable.} The Phase-0/1 v1 retraction, the adaptive-v1 development negative, the falsified predictions P1 and P6, the confounded X2 arm, and the corrected 34$\times$ ceiling claim are all preserved verbatim in the artifact; later results reinterpret none of them.
\item \textbf{The 7B confirmatory path is unchanged} and remains the only route to any real-model claim in this family.
\end{enumerate}

\section*{Code, materials and AI-use disclosure}
A public research artifact accompanies this series; every quantitative value in this report is read from frozen receipts in that artifact (the development-stress receipts, the attribution protocol and receipt, the three-tier ceiling bounds, the $\kappa$ freeze, the fresh-assurance and falsifiability-audit receipts, the licensing receipt, and the challenger development and assurance receipts), and the gate implementation with its known-answer battery is part of the artifact's source tree. Language models and coding agents were used for literature search, analysis, drafting and implementation assistance and are not authors, independent reviewers or human experts. Same-context analysis is labeled as such throughout; no claim in this paper has received independent review, and none is represented as having done so.

\begin{thebibliography}{99}
\sloppy

\bibitem{chen2023skillit}
M.~F.~Chen, N.~Roberts, K.~Bhatia, J.~Wang, C.~Zhang, F.~Sala, and C.~R\'e,
``Skill-It! A Data-Driven Skills Framework for Understanding and Training Language Models,''
\emph{Advances in Neural Information Processing Systems} (2023), \emph{arXiv:2307.14430}.

\bibitem{xie2023doremi}
S.~M.~Xie, H.~Pham, X.~Dong, N.~Du, H.~Liu, Y.~Lu, P.~Liang, Q.~V.~Le, T.~Ma, and A.~W.~Yu,
``DoReMi: Optimizing Data Mixtures Speeds Up Language Model Pretraining,''
\emph{arXiv:2305.10429} (2023).

\bibitem{tirumala2023d4}
K.~Tirumala, D.~Simig, A.~Aghajanyan, and A.~S.~Morcos,
``D4: Improving LLM Pretraining via Document De-Duplication and Diversification,''
\emph{Advances in Neural Information Processing Systems} (2023), \emph{arXiv:2308.12284}.

\bibitem{abbas2023semdedup}
A.~Abbas, K.~Tirumala, D.~Simig, S.~Ganguli, and A.~S.~Morcos,
``SemDeDup: Data-Efficient Learning at Web-Scale through Semantic Deduplication,''
\emph{arXiv:2303.09540} (2023).

\bibitem{lin2024rho1}
Z.~Lin, Z.~Gou, Y.~Gong, X.~Liu, Y.~Shen, R.~Xu, C.~Lin, Y.~Yang, J.~Jiao, N.~Duan, and W.~Chen,
``Rho-1: Not All Tokens Are What You Need,''
\emph{arXiv:2404.07965} (2024).

\bibitem{xia2024less}
M.~Xia, S.~Malladi, S.~Gururangan, S.~Arora, and D.~Chen,
``LESS: Selecting Influential Data for Targeted Instruction Tuning,''
\emph{arXiv:2402.04333} (2024).

\bibitem{yu2024mates}
Z.~Yu, S.~Das, and C.~Xiong,
``MATES: Model-Aware Data Selection for Efficient Pretraining with Data Influence Models,''
\emph{arXiv:2406.06046} (2024).

\bibitem{li2025mass}
J.~Li et al.,
``MASS: Mathematical Data Selection via Skill Graphs for Pretraining Large Language Models,''
\emph{International Conference on Machine Learning} (2025), \emph{arXiv:2503.14917}.

\bibitem{he2026stat}
Y.~He et al.,
``STAT: Skill-Targeted Adaptive Training,''
\emph{International Conference on Learning Representations} (2026), \emph{arXiv:2510.10023}.

\bibitem{wei2026steps}
Y.~Wei, L.~Du, X.~Yu, Y.~Feng, and A.~Li,
``Towards Compositional Generalization of LLMs via Skill Taxonomy Guided Data Synthesis (STEPS),''
\emph{arXiv:2601.03676} (2026).

\bibitem{chen2026structural}
M.~C.~Chen, M.~Miller, B.~Sch\"olkopf, and S.~Guo,
``On the Emergence and Test-Time Use of Structural Information in Large Language Models,''
\emph{arXiv:2601.17869} (2026).

\bibitem{ling2026rex}
Z.~Ling et al.,
``Reusable Experiences: Latent Routing and Modular Composition in LLMs,''
\emph{Annual Meeting of the Association for Computational Linguistics} (2026).

\bibitem{butler2010}
A.~C.~Butler,
``Repeated Testing Produces Superior Transfer of Learning Relative to Repeated Studying,''
\emph{Journal of Experimental Psychology: Learning, Memory, and Cognition} \textbf{36}(5), 1118--1133 (2010).

\bibitem{yiu2026verified}
S.~C.~Yiu,
``\RAKL{} research series: method-evolution and verified-discovery source packages,''
\RAKL{} research artifact (2026).

\end{thebibliography}
\end{document}
